\chapter{\texorpdfstring{\(L^p\)}{Lp} 空间的包含关系}

\begin{solution}{finite-lp-inclusion-sharp}
因 \(0<\mu(X)<\infty\)，可取非零常值函数
\(f=c\mathbf1_X\)。则
\[
 \|f\|_p=|c|\mu(X)^{1/p},\qquad
 \|f\|_q=|c|\mu(X)^{1/q},
\]
两边恰好取等。因此任何对所有函数成立的常数不能小于
\(\mu(X)^{1/p-1/q}\)。
若 \(q=\infty\)，则 \(\|f\|_\infty=|c|\)，仍由同一常函数取等；
这里使用 \(1/\infty=0\)。
\end{solution}

\begin{solution}{finite-lp-strict}
先设 \(q<\infty\)。取 \(f(x)=x^{-\alpha}\)，其中
\(1/q\le\alpha<1/p\)。因为
\[
 \int_0^1x^{-\alpha r}\,dx<\infty
\quad\Longleftrightarrow\quad \alpha r<1,
\]
故 \(f\in L^p(0,1)\)，而 \(f\notin L^q(0,1)\)。
若 \(q=\infty\)，取 \(0<\alpha<1/p\)。其 \(p\) 次积分为
\(1/(1-\alpha p)<\infty\)，而对每个 \(M>1\)，
\(\{f>M\}=(0,M^{-1/\alpha})\) 有正测度，故不属于 \(L^\infty\)。
\end{solution}

\begin{solution}{infinite-no-inclusion}
下面的函数只在指定的正半轴区间内计算幂，其他点取零。
取 \(1/q<\alpha\le1/p\)，则
\[
 f(x)=x^{-\alpha}\mathbf1_{(1,\infty)}
\]
满足 \(f\in L^q\setminus L^p\)。取
\(1/q\le\beta<1/p\)，则
\[
 g(x)=x^{-\beta}\mathbf1_{(0,1)}
\]
满足 \(g\in L^p\setminus L^q\)。条件分别来自
\(\int_1^\infty x^{-s}\) 在 \(s>1\) 时收敛和
\(\int_0^1x^{-s}\) 在 \(s<1\) 时收敛。
\end{solution}

\begin{solution}{ellp-inclusion-strict}
正文定理给 \(\ell^p\subset\ell^q\)。令
\(x_n=(n+1)^{-\alpha}\)，其中
\(1/q<\alpha\le1/p\)，则
\(\sum|x_n|^q<\infty\) 而 \(\sum|x_n|^p=\infty\)，故包含严格。
有限支撑序列属于每个 \(0<p\le\infty\) 的 \(\ell^p\)。
\end{solution}

\begin{solution}{ellp-tends-zero}
若不趋于零，则存在 \(\varepsilon>0\) 和无限多个 \(n\) 使
\(|x_n|\ge\varepsilon\)，从而
\(\sum|x_n|^p\ge\sum\varepsilon^p=\infty\)，矛盾。并且对每个 \(n\)，
\[
 |x_n|^p\le\sum_k|x_k|^p,
\]
故 \(\|x\|_\infty\le\|x\|_p\)。
\end{solution}

\begin{solution}{essential-sup-examples}
Dirichlet 函数只在可数集上取 \(1\)，故本质上等于零，
\(\|f\|_\infty=0\)。Thomae 函数也只在可数集上非零，所以本质上确界
为零。函数 \(x^{-1/2}\mathbf1_{(0,1)}\) 对每个 \(M>0\) 都在
\((0,\min\{1,M^{-2}\})\) 上大于 \(M\)，该集合正测，故本质上确界为
\(+\infty\)。
\end{solution}

\begin{solution}{linfty-completeness-direct}
若 \((f_n)\) 是 Cauchy 列，给定 \(k\)，取 \(N_k\) 使
\[
 m,n\ge N_k\Longrightarrow
 \|f_m-f_n\|_\infty<2^{-k}.
\]
各 \(k\ge1\)、\(m,n\ge N_k\) 对应的集合
\(\{|f_m-f_n|\ge2^{-k}\}\) 都是可测零集：因本质范数严格小于此阈值，
可先选一个更小的几乎处处上界。将这些集合及
\(\{|f_{N_1}|>\|f_{N_1}\|_\infty\}\) 作可数并，记为 \(Z\)。
在 \(X\setminus Z\) 上整列一致 Cauchy，故逐点存在有限极限。
定义该极限为 \(f\)，在 \(Z\) 上置零。它是
\(f_n\mathbf1_{X\setminus Z}\) 的点态极限，故可测，且满足
\[
 \|f_n-f\|_\infty\le2^{-k}\quad(n\ge N_k).
\]
取 \(k=1\) 得 \(|f|\le\|f_{N_1}\|_\infty+1/2\) 几乎处处，
故 \(f\in L^\infty\)。再令 \(k\to\infty\)，得到原列在本质上确界范数下
收敛。整个论证只在可测零集上补零，不要求测度空间完备。
\end{solution}

\begin{solution}{lp-limit-characteristic}
\[
 \|\mathbf1_E\|_p
 =\left(\int_E1\,d\mu\right)^{1/p}
 =\mu(E)^{1/p}.
\]
若 \(0<\mu(E)<\infty\)，该值趋于 \(1=\|\mathbf1_E\|_\infty\)；若
\(\mu(E)=0\)，所有范数为零；若 \(\mu(E)=\infty\)，指示函数不属于任何
有限 \(L^p\)，但属于 \(L^\infty\)。
\end{solution}

\begin{solution}{lp-limit-unbounded}
在 \((0,1)\) 上取 \(f(x)=|\log x|\)。先在
\([e^{-R},e^{-\delta}]\)（\(0<\delta<R\)）上使用经典积分换元
\(u=-\log x\)，再令 \(R\uparrow\infty,\delta\downarrow0\)，
由单调收敛及连续函数的两种积分一致性得到
\[
 \int_0^1|\log x|^p\,dx=\int_0^\infty u^pe^{-u}\,du<\infty
\]
对每个 \(0<p<\infty\) 成立：原点附近 \(u^p\) 可积，无穷远处
\(u^pe^{-u}\le e^{-u/2}\) 最终成立。
对每个 \(M>0\)，\(\{f>M\}=(0,e^{-M})\) 的测度为 \(e^{-M}>0\)，
所以 \(f\notin L^\infty\)。若 \(\|f\|_p\) 不趋于无穷，则存在趋于无穷的指数序列上
范数有界。因为 \((0,1)\) 的测度为 \(1\)，有限测度包含估计说明
\(p\mapsto\|f\|_p\) 单调不减；所以该子列的有界性推出所有充分大的 \(p\)
上的范数一致有界。定理 \ref{thm:lp-limit-linfty} 的第二部分于是推出
\(f\in L^\infty\)，矛盾。
\end{solution}

\begin{solution}{local-lp-singularity}
先在环域 \(\rho<|x|<R\)（\(0<\rho<R<\infty\)）上计算。
被积函数连续有界，环域边界零测。把函数在包含环域的矩形内延拓为零，
用矩形分割的下、上阶梯函数夹住它，其 Lebesgue 积分分别就是 Darboux
下、上和；两和之差趋零，便得到这里两种积分的一致性。
因此可以先用经典多重积分的极坐标换元，得到
\[
 \int_{\rho<|x|<R}|x|^{-\alpha p}\,dx
 =C_n\int_\rho^R r^{n-1-\alpha p}\,dr,
\]
其中 \(C_n>0\) 是单位球面的面积（\(n=1\) 时 \(C_1=2\)，
公式直接由正负两个区间相加得到）。取 \(R=1\) 并令 \(\rho\downarrow0\)，
使用非负函数的单调收敛，得
\[
 \int_{|x|<1}|x|^{-\alpha p}\,dx
 =C_n\int_0^1r^{n-1-\alpha p}\,dr,
\]
当且仅当 \(\alpha p<n\) 收敛。因此
\(|x|^{-\alpha}\in L^p_{\rm loc}\) 恰在 \(\alpha p<n\)（若只在
\(\mathbb R^n\setminus\{0\}\) 局部则无此限制）。另取 \(\rho=1\) 并令
\(R\uparrow\infty\)，同一单调收敛论证给出无穷远处积分
\(\int_1^\infty r^{n-1-\alpha p}dr\) 当且仅当 \(\alpha p>n\) 收敛。
\end{solution}

\begin{solution}{log-convex-lp-norms}
若 \(\theta=0\) 或 \(1\)，结论为端点恒等式。以下设 \(0<\theta<1\)。
若某个端点范数为零，\(f=0\) 几乎处处，结论成立。其余情形端点范数为正。
先设 \(p_1<\infty\)，令
\[
 a=\frac{(1-\theta)p_\theta}{p_0},\qquad
 b=\frac{\theta p_\theta}{p_1};
\]
由指数关系 \(a+b=1\)。写
\[
 |f|^{p_\theta}
 =(|f|^{p_0})^a(|f|^{p_1})^b.
\]
对共轭指数 \(1/a,1/b\) 使用 Hölder：
\[
 \int|f|^{p_\theta}
 \le\left(\int|f|^{p_0}\right)^a
     \left(\int|f|^{p_1}\right)^b.
\]
取 \(p_\theta\) 次根即得所求。
若 \(p_1=\infty\)，则 \(p_\theta=p_0/(1-\theta)<\infty\)，故
\[
 \int|f|^{p_\theta}
 \le\|f\|_\infty^{p_\theta-p_0}\int|f|^{p_0}.
\]
取 \(p_\theta\) 次根并用 \(p_0/p_\theta=1-\theta\)，即得
\(\|f\|_{p_\theta}\le\|f\|_{p_0}^{1-\theta}\|f\|_\infty^\theta\)。
两种证明中的右侧均有限，也同时确认了 \(f\in L^{p_\theta}\)。
\end{solution}

\begin{solution}{lp-membership-phase-diagram}
原点附近 \(1+|x|\asymp1\)，所以需要
\[
 \alpha p<n.
\]
无穷远处
\[
 |x|^{-\alpha}(1+|x|)^{-\beta}
 \asymp |x|^{-(\alpha+\beta)},
\]
所以需要
\[
 (\alpha+\beta)p>n.
\]
两条严格不等式同时成立时属于 \(L^p(\mathbb R^n)\)；等号情形产生
\(\int dr/r\) 的对数发散。若 \(\alpha\) 或 \(\beta\) 为负，同一渐近
计算仍适用。

固定 \(p>0\)，记 \(a=n/p\)。可积参数恰为开区域
\[
 \{(\alpha,\beta):\alpha<a,\ \beta>a-\alpha\}.
\]
它位于竖线 \(\alpha=a\) 左侧、斜线 \(\beta=a-\alpha\) 上方；
两条边界均不包含。下图画出该区域在坐标窗口中的部分，阴影向上及左上方延伸。

\begin{center}
\begin{tikzpicture}[x=1cm,y=1cm]
 \fill[gray!18] (-1.5,3.5)--(1.5,3.5)--(1.5,0)--(-1.5,3)--cycle;
 \draw[->] (-1.8,0)--(3.2,0) node[right] {\(\alpha\)};
 \draw[->] (0,-0.8)--(0,3.8) node[above] {\(\beta\)};
 \draw[dashed] (1.5,-0.6)--(1.5,3.5);
 \draw[dashed] (-1.5,3)--(2.1,-0.6);
 \node[below left] at (0,0) {\(0\)};
 \node[below right] at (1.5,0) {\(a\)};
 \node[left] at (0,1.5) {\(a\)};
 \node at (0.55,2.9) {可积};
 \node[right] at (1.5,3.25) {\(\alpha=a\)};
 \node[right] at (2.05,-0.45) {\(\beta=a-\alpha\)};
\end{tikzpicture}
\end{center}
这里 \(a>0\) 只是示意比例；对所有 \(p>0\) 取这些截面，即得到题目所求的
三参数区域。原点取零不会改变积分。
\end{solution}
